\chapter{Приведение квадратичных форм в линейном пространстве к каноническому виду. Положительно определенные квадратичные формы. Критерий Сильвестра.}\label{chapter25}

\section[Билинейные и квадратичные формы.]{Билинейные и квадратичные формы\footnote{Рекомендую ознакомиться с написанными самим Чубаровым И.А. материалами по этой главе по этой ссылке: \href{https://drive.google.com/drive/u/0/folders/0BzuzEyNkpwYDcFhhV1l2N1lhY2s}{$https://drive.google.com/drive/...$}}.}
Пусть $L$ "--- линейное пространство над полем $K$. При изложении вопроса достаточно считать поле скаляров $K$ множеством действительных чисел~$\bbR$.
\begin{defn}  
Функция $b(x,y)\colon L\times L\rightarrow K$ называется \textit{билинейной функцией}, если она линейна по каждому аргументу, то есть
\begin{enumerate}
\item 
$b(\alpha_1 x_1+\alpha_2 x_2,y)=\alpha_1 b(x_1,y)+\alpha_2 b(x_2,y)$ $\forall x_1,x_2,y \in L, \alpha_1,\alpha_2 \in K$,
\item 
$b(x,\beta_1 y_1+\beta_2 y_2)=\beta_1 b(x,y_1)+\beta_2 b(x,y_2)$  $\forall x,y_1,y_2 \in L, \beta_1,\beta_2 \in K$.
\end{enumerate}
\end{defn}

Пусть $\dim L=n$, $\overline{e}=\norm{e_1 \ldots e_n}$ "--- базис $L$. Обозначим $b_{ij}=b(e_i,e_j), 1\le i,j \le n$.
\begin{defn}
Матрицу $B=B_e=(b_{ij})$ называют \textit{матрицей билинейной функции} $B$ в базисе $e_1$, \ldots,~$e_n$.
\end{defn}
  
Получим координатную запись. Пусть $x=\sum\limits_{i=1}^nx_ie_i,y=\sum\limits_{j=1}^ny_je_j$, тогда
\begin{equation}\label{24.1.bilform}
b(x,y)=b(\sum_{i=1}^nx_ie_i,\sum_{j=1}^ny_je_j)=\sum_{i,j=1}^nx_iy_jb(e_i,e_j)=\sum_{i,j=1}^nx_ib_{ij}y_j=X^TBY,
\end{equation}
где $X=\begin{pmatrix}
x_1 \\ \vdots \\ x_n
\end{pmatrix} Y=\begin{pmatrix}
y_1 \\ \vdots \\ y_n
\end{pmatrix}$.
\begin{defn}
Запись билинейной функции в виде многочлена~\eqref{24.1.bilform} называют \textit{билинейной формой}.
\end{defn}

\begin{notion} 
По традиции билинейной формой называют и саму функцию.
\end{notion}

\begin{stt}
Пусть $e=(e_1,...,e_n), e'=(e'_1,...,e'_n)$ "--- два разных базиса пространства $L$, $S=S_{e\rightarrow e'}$ "--- матрица перехода от базиса $e$ к базису $e'$, а $B$ и $B'$ "--- матрицы билинейной формы $b$ в базисах $e$ и $e'$ соответственно. Тогда 
\begin{equation}\label{24.1.transformation}
B'=S^TBS
\end{equation}

\end{stt}
Из формулы \eqref{24.1.transformation} следует, что ранг матрицы $B$ и знак её определителя (если он не равен нулю) не зависят от выбора базиса.
\begin{defn}
Билинейная форма $b(x,y)$ называется \textit{симметрической}, если $\forall x,y \in L \hookrightarrow b(x,y)=b(y,x)$.
\end{defn}
\begin{stt}
Матрица симметрической билинейной формы в любом базисе является симметрической, т.е.~$B^T=B$.
\end{stt}

\begin{defn}
\textit{Квадратичной функцией (формой)}, порождённой симметрической билинейной формой $ b(x,y) $, называется функция $k(x)=b(x,x)$  $\forall x \in L$.
\end{defn}
\begin{stt}
Для любой квадратичной функции $k(x)$ существует единственная симметрическая билинейная форма $b(x,y)$, такая, что $k(x)=b(x,x), \forall x \in L$
\end{stt}
\begin{proof}
$k(x+y)=b(x+y,x+y)=b(x,x)+2b(x,y)+b(y,y)=k(x)+k(y)+2b(x,y) \Rightarrow b(x,y)=\dfrac{k(x+y)-k(x)-k(y)}{2}$.
\end{proof}
\begin{defn}
\textit{Матрицей квадратичной формы} называют матрицу породившей её симметричной билинейной формы.
\end{defn}
Квадратичная форма записываема в виде 			  \begin{equation}\label{24.1.common}
k(x)=b_{11}x_1^2+\ldots+b_{nn}x_n^2+2\sum\limits_{1\le i<j\le n} b_{ij}x_ix_j.
\end{equation} 
\section{Приведение квадратичных форм в линейном пространстве к каноническому виду}
\begin{defn}
Квадратичная форма вида $q(x)=\sum\limits_{i=1}^n\alpha_ix_i^2$ называется \textit{диагональной}. Она называется \textit{канонической}, если $\alpha_i=\pm 1;0$ и 
\begin{equation}\label{24.2.canonical}
q(x)=\sum\limits_{i=1}^py_i^2-\sum\limits_{i=p+1}^{p+q}y_i^2
\end{equation} 
Числа $p$ и $q$ называются \textit{положительным} и \textit{отрицательным индексами инерции}, $\sigma=p-q$ "--- \textit{сигнатурой}.
\end{defn}
\begin{thm}[о приведении квадратичной формы к каноническому виду]
Для любой квадратичной формы~\eqref{24.1.common} существует такая невырожденая $(\det S \neq 0)$ замена переменных $X=SY$, что в новых переменных она принимает канонический вид~\eqref{24.2.canonical}.
\end{thm}
\begin{proof} Воспользуемся \textit{алгоритмом Лагранжа} выделения полных квадратов.\\
(1) Допустим,	что $\exists i$: $b_{ii} \neq 0$. При 	необходимости перенумеровав переменные, можем считать, что $b_{11} \neq 0$. Тогда выделим в квадратичной форме все одночлены, содержащие $x_1$, и дополним это выражение до квадрата:
\begin{equation*}\begin{array}{crl}
k(x)=b_{11}x_1^2+2\sum\limits_{j=2}^nb_{1j}x_1x_j+(\sum\limits_{i=2}^nb_{ii}x_i^2+2\sum\limits_{2\le i<j\le n} b_{ij}x_ix_j)=\\
=b_{11}(x_1+\sum\limits_{j=2}^n\dfrac{b_{1j}}{b_{11}}x_j)^2+(\sum\limits_{i=2}^nb_{ii}x_i^2+2\sum\limits_{2\le i<j\le n} b_{ij}x_ix_j-(\sum\limits_{j=2}^n\dfrac{b_{1j}}{b_{11}}x_j)^2)=\\
=b_{11}(x_1+\sum\limits_{j=2}^n\dfrac{b_{1j}}{b_{11}}x_j)^2+k_1(x_2,\ldots, x_n)
\end{array}\end{equation*}
Тогда сделаем замену $\widetilde{x_1}=x_1+\sum\limits_{j=2}^n\dfrac{b_{1j}}{b_{11}}x_j, \widetilde{x_2}=x_2,\ldots \widetilde{x_n}=x_n$. Квадратичная форма $k_1(\widetilde{x_2},\ldots, \widetilde{x_n})$ не зависит от $\widetilde{x_1}$, и к ней можно применить тот же метод, и так далее, в результате получится квадратичная форма $\sum\limits_{i=1}^r\alpha_i\widetilde{x_i}^2 (\alpha_1\alpha_2\ldots\alpha_r \neq 0, r=\rg B)$. Остаётся сделать замену $y_i=\sqrt{\abs{\alpha_i}}\widetilde{x_i}, i=\overline{1,r}; y_k=\widetilde{x_k}, k=\overline{r+1,n}$.\\
(2) Препятствие к выделению квадратов может возникнуть, если на каком-то этапе получена форма, все диагональные элементы матрицы которой равны нулю, но есть ненулевые элементы вне диагонали. Перенумеровав при необходимости переменные, добьёмся $b_{12} \neq 0$. Тогда сделаем подготовительную замену $x_1=x'_1-x'_2,x_2=x'_1+x'_2,x_j=x'_j (j \ge 3)$:
\begin{equation*}
k(x')=2b_{12}({x'}_1^2-{x'}_2^2)+q'(x'),
\end{equation*}
причём форма $q'(x')$ не содержит с члена с ${x'}_1^2$. Теперь условие пункта 1 выполняется.
\end{proof}

\begin{thm}\label{24.2.thm2}
Пусть невырожденная замена $X=SY$ приводит квадратичную форму $k(x)$ к каноническому виду \eqref{24.2.canonical}. Если другая невырожденная замена $X=TZ$ приводит форму к каноническому виду $\widetilde q(x)=\sum\limits_{i=1}^sz_i^2-\sum\limits_{i=s+1}^{s+t}z_i^2$, то $p=s,q=t,$ причём $p+q=\rg B$
\end{thm} 
Теорему \ref{24.2.thm2} оставим без доказательства, только заметим, что равенство $p+q=\rg B$ следует из сохранения ранга матрицы $B$ при замене базиса.

\section[Положительно определенные квадратичные формы.]{Положительно определенные квадратичные формы\footnote{Рекомендую ознакомиться с написанными самим Чубаровым И.А. материалами по этому билету по этой ссылке: \href{https://drive.google.com/drive/u/0/folders/0BzuzEyNkpwYDcFhhV1l2N1lhY2s}{$https://drive.google.com/drive/...$}}. }
  \begin{defn}
  Квадратичная функция $k(x)$ на линейном пространстве~$L$ (здесь пространство задаётся над полем вещественных чисел: $K=\bbR$) называется \textit{положительно определённой}, если $\forall x \in L$: $x \neq 0 \hookrightarrow k(x)>0;$ \textit{отрицательно определённой}, если $\forall x \in L$: $x \neq 0 \hookrightarrow k(x)<0;$ \\
\textit{неотрицательно определённой (положительно полуопределённой)}, если $\forall x \in L \hookrightarrow k(x)\ge 0;$ \\
\textit{неположительно определённой (отрицательно полуопределённой)}, если $\forall x \in L \hookrightarrow k(x)\le 0;$ \\
\textit{знаконеопределённой}, если $\exists x_1,x_2 \in L$ и $k(x_1)>0$, $k(x_2)<0.$ \\
  \end{defn}
  Пусть в некотором базисе квадратичная функция записана в виде квадратичной формы
  \begin{equation}\label{25.1.common}
  k(x)=b_{11}x_1^2+\ldots+b_{nn}x_n^2+2\sum\limits_{1\le i<j\le n} b_{ij}x_ix_j.
  \end{equation}
с матрицей $B=(b_{ij})$.
  \begin{lemm}\label{25.1.lemm}
  Квадратичная форма тогда и только тогда является положительно определенной, когда она приводится к диагональному виду $\sum\limits_{i=1}^n\alpha_iz_i^2, \alpha_i>0, i=\overline{1,n}$ или, что равносильно, каноническому виду $\sum\limits_{i=1}^ny_i^2$.
  \end{lemm}
  \begin{notion}
  От диагонального вида можно прийти к каноническому при помощи замены $y_i=\sqrt{\alpha_i}z_i$.
  \end{notion}
  \begin{proof}
  То, что диагональная форма со всеми положительными коэффициентами положительно определена, очевидно. \\
  Обратно, допустим, что данная положительно определённая форма приводится к виду $\sum\limits_{i=1}^py_i^2-\sum\limits_{i=p+1}^{p+q}y_i^2$. Если, вопреки доказываемому, $p<n$, то $k(0,0,...,1)\le 0$. Получаем противоречие с условием.
  \end{proof}
\section{Критерий Сильвестра}
  \begin{thm} [Критерий Сильвестра\rindex{критерий!Сильвестра}]
  Для положительной определенности квадратичной формы $k(x)$ в $\bbR^n$ необходимо и достаточно, чтобы все \underline{главные миноры} её матрицы $B$, имеющие вид
  \begin{equation}
  \Delta_m=\det \begin{Vmatrix}
  b_{11} & b_{12} & \cdots & b_{1m} \\
  b_{21} & b_{22} & \cdots & b_{2m} \\
  \vdots & \vdots & \ddots & \vdots \\
  b_{m1} & b_{m2} & \cdots & b_{mm} \\
  \end{Vmatrix}(b_{ij}=b_{ji} \forall i,j), m=\overline{1,n},
  \end{equation}
были положительными
  \end{thm}
  \begin{proof} \leavevmode
  \begin{itemize}[wide, labelwidth=!, labelindent=\parindent]
  \item[\fbox{дост.:}] Пусть дано, что все главные миноры матрицы квадратичной формы положительны, и надо доказать, что она является положительно определенной. Воспользуемся методом математической индукции и леммой~\ref{25.1.lemm}. 

  Случай для \underline{$n=1$} очевиден. 
  
  Допустим, что \underline{$n>1$} и из положительности главных миноров матрицы квадратичной формы вплоть до $(n-1)$-го порядка включительно следует возможность приведения квадратичной формы от $n-1$ переменных $x_1,...,x_{n-1}$ к виду $k(x)=\sum\limits_{i=1}^{n-1}x_i^2$. Покажем, что достаточность имеет место и в случае $n$ переменных.
  
  В выражении квадратичной формы от переменных $x_1,...,x_n$ выделим слагаемые, содержащие $x_n$:
  \begin{equation*}
  k(x)=\sum_{j=1}^{n-1}\sum_{i=1}^{n-1}b_{ij}x_jx_i+2\sum_{j=1}^{n-1}b_{jn}x_jx_n+b_{nn}x_n^2.
  \end{equation*}
  Двойная сумма $\sum\limits_{j=1}^{n-1}\sum\limits_{i=1}^{n-1}b_{ij}x_jx_i=k^*(x)$ есть квадратичная форма, зависящая от $n-1$ переменной, причем главные миноры её матрицы совпадают с главными минорами $k(x)$ до порядка $n-1$ включительно, которые, по условию, положительны. Отсюда следует, по предположению индукции, что квадратичная форма $k^*(x)$ положительно определена и для нее существует невырожденная замена переменных $x_j=\sum\limits_{i=1}^{n-1}\sigma_{ji}y_i, j=\overline{1,n-1}$, приводящая ее к каноническому виду $k^*(x)=\widetilde k^*(y)=\sum\limits_{i=1}^{n-1}y_i^2$.
  
  Запишем квадратичную форму в новых переменных:
  \begin{equation*}
  k(x)=\widetilde k(y_1,...,y_{n-1},x_n)=\sum\limits_{i=1}^{n-1}y_i^2+2\sum_{j=1}^{n-1}b'_{jn}y_jx_n+b_{nn}x_n^2
  \end{equation*}
и выделим полные квадраты по $y_1,...,y_{n-1}$:
  \begin{equation*}\begin{array}{crl}
  k=\sum\limits_{i=1}^{n-1}(y_i^2+2b'_{in}y_ix_n+{b'_{in}}^2x_n^2)+(b_{nn}-\sum\limits_{i=1}^{n-1}{b'_{in}}^2)x_n^2 =\\
  =\sum\limits_{i=1}^{n-1}z^2+b''_{nn}x_n^2,
  \end{array}\end{equation*}
где введено обозначение $ b''_{nn}=b_{nn}-\sum\limits_{i=1}^{n-1}{b'_{in}}^2 $ и произведена замена $z_i=y_i+b'_{in}x_n, i=\overline{1,n-1}, x_n=x_n$. Эта замена, очевидно, невырожденная.

  Теперь вспомним, что определитель матрицы квадратичной формы сохраняет знак при замене базиса. По условию определитель матрицы $B$ квадратичной функции в исходном базисе положительный, поскольку является главным минором порядка $n$. Но из выражения для $k(x)$ в конечном базисе мы получаем, что определитель матрицы квадратичной формы $k$ равен $b''_{nn}$. Поэтому $b''_{nn}>0$ и можно ввести переменную $z_n=\sqrt{b''_{nn}}x_n$, в результате чего получаем канонический вид: $k=\sum\limits_{i=1}^nz_i^2$.
  
  \item[\fbox{необх.:}] Дано, что квадратичная функция положительно определена, и надо доказать положительность главных миноров её матрицы. Снова применим индукцию по числу переменных $n$. 
  
  Для \underline{n=1} это ясно.
  
  Пусть \underline{n>1} и для форм от меньшего числа переменных утверждение теоремы верно. Поскольку квадратичная форма $k^*(x)=\sum\limits_{j=1}^{n-1}\sum\limits_{i=1}^{n-1}b_{ij}x_jx_i$ является положительно определенной (её значения –- это значения $k(x)$ при $x_n=0$), то по предположению индукции ее главные миноры, совпадающие с главными минорами матрицы $B$ до порядка ${n-1}$, положительны. А определитель самой матрицы $B$, который является главным минором порядка $n$, положителен, поскольку $k(x)$ приводится к каноническому виду $k=\sum\limits_{i=1}^nz_i^2$, и определитель матрицы полученной при этом квадратичной формы равен 1 и имеет такой же знак, как и определитель матрицы $B$.
  \end{itemize}
  \end{proof}
  
  \begin{cons} [Критерий Сильвестра для отрицательной определенности]
  Для отрицательной определенности квадратичной формы  $k(x)$ в $\bbR^n$ необходимо и достаточно, чтобы все главные миноры ее матрицы $B$ имели чередующиеся знаки, начиная с минуса, т.е. $(-1)^m\Delta_m>0,m=\overline{1,n}$.
  \end{cons}
  \begin{proof}
  Рассмотрим форму $-k(x)$ c матрицей $B'=-B=(-b_{ij})$: ее положительной определённости, по критерию Сильвестра, равносильно условие
  \begin{equation}
  \Delta_m=\det \begin{Vmatrix}
  -b_{11} & -b_{12} & \cdots & -b_{1m} \\
  -b_{21} & -b_{22} & \cdots & -b_{2m} \\
  \vdots  & \vdots  & \ddots & \vdots \\
  -b_{m1} & -b_{m2} & \cdots & -b_{mm} \\
  \end{Vmatrix}=(-1)^m\Delta_m>0, m=\overline{1,n}.
  \end{equation}
  \end{proof}